% TOP_PROBLEM: {"id": 2913, "title": "Kirby Problem 4.37", "category_id": 7, "category": {"id": 7, "name": "topology", "display_name": "Topology", "description": "Properties preserved under continuous deformations.", "slug": "topology", "order_index": 7, "created_at": "2026-07-31T15:26:25.671Z"}, "status": "open"}
% FIRST_POSTED: null
% ENTRY_KIND: statement_only
% Imported from: batch13.tex; UnsolvedMath ID 2913
\documentclass[11pt]{article}
\usepackage[margin=25mm]{geometry}
\usepackage{fontspec}
\setmainfont{Latin Modern Roman}
\usepackage{amsmath,amssymb,mathtools}
\usepackage{unicode-math}
\setmathfont{Latin Modern Math}
\usepackage{xurl,hyperref}
\hypersetup{colorlinks=true,urlcolor=blue}
\setcounter{secnumdepth}{0}
\setlength{\parindent}{0pt}
\setlength{\parskip}{5pt}
\setlength{\emergencystretch}{3em}
\newcommand{\sref}[2]{\hyperlink{source:#1:#2}{[S#2]}}
\newcommand{\cataloguescope}[1]{\paragraph{Recorded target.}#1}
\begin{document}
% UnsolvedMath ID 2913; source catalogue code KP-4.37
\setcounter{section}{2912}
\section{Kirby Problem 4.37}\label{2913}
{\small\sffamily UnsolvedMath ID 2913 · Topology}
\subsection{Definitions and mathematical statement}

\paragraph{Shared notation.}$\mathbb N=\{1,2,\ldots\}$, and $\mathbb Z,\mathbb Q,\mathbb R,\mathbb C$ denote the integers, rationals, reals and complex numbers. Any other conventions are specified locally.


A projective plane in $S^4$ is an embedding of $\mathbb{RP}^2$. A two-knot is an embedding $S^2\hookrightarrow S^4$, with no assumption that it is unknotted. Connected sum of embedded surfaces is formed by removing a small standard disk from each in separate four-balls and joining the boundary circles by a tube in the connected sum $S^4\#S^4\cong S^4$.

For orientable genus $g\ge0$, the standard unknotted surface is the boundary of a standard genus-$g$ handlebody in an equatorial $S^3\subset S^4$. For the nonorientable models, write $P_+$ and $P_-$ for the standard unknotted projective planes with normal Euler numbers $+2$ and $-2$. One explicit model is the normalized Veronese embedding
\[
\mathbb{RP}^2\longrightarrow S^4\subset\operatorname{Sym}_0(3,\mathbb R),\qquad
[v]\longmapsto\sqrt{\frac32}\left(vv^{\mathsf T}-\frac13I_3\right),\quad |v|=1,
\]
where $\operatorname{Sym}_0(3,\mathbb R)$ is the five-dimensional space of real symmetric trace-zero matrices with its Euclidean matrix norm; reflection gives the other sign. The standard nonorientable genus-$h$ surfaces are connected sums of $h$ copies of $P_+$ or $P_-$, formed by joining standard local disks by tubes. Their normal Euler number is the sum of the signs $\pm2$. An embedding is topologically or smoothly unknotted if it is ambiently isotopic, in that category, to a corresponding standard model.
\paragraph{Statement recorded in the supplied source.}
In each of the smooth and the locally flat categories, is every embedded projective plane $P\subset S^4$ ambiently isotopic to
\[
S\#P_+\quad\text{or}\quad S\#P_-
\]
for some two-knot $S\subset S^4$ in the same category? This is the Kinoshita decomposition question; the sphere summand is allowed to be knotted. The source also records these subsidiary questions:
\begin{enumerate}
\item[(i)] Does there exist an embedded Klein bottle $\Sigma\subset S^4$ with $|e(\Sigma)|=4$ that has no unknotted projective-plane connected-sum summand? The normal Euler number $e(\Sigma)$ is the Euler number of its normal two-plane bundle, with local orientation coefficients.
\item[(ii)] For a projective plane $P\subset S^4$ and a tubular neighborhood $\nu P$, must the kernel of the inclusion-induced map
\[
\pi_1\bigl(\partial(S^4\setminus\operatorname{int}\nu P)\bigr)
\longrightarrow\pi_1(S^4\setminus P)
\]
have order four? Use a boundary basepoint and its image; changing it by a path only conjugates the kernel.
\end{enumerate}
As with the main decomposition question, these embedding questions can be considered in the smooth or locally flat category. \sref{2913}{2}
\subsection{Short English statement}
Ask whether every projective plane splits off an unknotted projective plane, whether an Euler-four Klein bottle can avoid such a summand, and whether each projective-plane peripheral inclusion has kernel of order four.
\subsection{Sources}
\begin{description}
\item[\hypertarget{source:2913:1}{S1}] ULAM.AI and the UnsolvedMath Contributors, UnsolvedMath: A Curated Collection of Open Mathematics Problems, record 2913 (KP-4.37). CC BY 4.0. \url{https://huggingface.co/datasets/ulamai/UnsolvedMath}.
\item[\hypertarget{source:2913:2}{S2}] R. İnanç Baykur, Robion C. Kirby and Daniel Ruberman, K3—A New Problem List in Low-Dimensional Topology, Mathematical Surveys and Monographs 295 (2026), author version, Problem 4.37 and Remarks, PDF page 220. \url{https://math.berkeley.edu/sites/default/files/surv-295-ruberman-watermarked-author-pdf.pdf}.
\end{description}
\label{end:2913}
\end{document}
